\documentclass{article}
\usepackage{apacite}
\usepackage{hyperref}
\usepackage{multicol}
\usepackage{natbib}
\usepackage{andi}
\usepackage[utf8]{inputenc}
\usepackage{array}
\usepackage[first=0,last=9]{lcg}
\setlength{\parindent}{0pt}
\usepackage{hyperref}
\usepackage{physics}
\usepackage{soul}
\usepackage{tikz-cd}
\usepackage{multirow}

\usepackage{courier} %% Sets font for listing as Courier.
\usepackage{listings, xcolor}
\lstset{
tabsize = 4, %% set tab space width
showstringspaces = false, %% prevent space marking in strings, string is defined as the text that is generally printed directly to the console
numbers = left, %% display line numbers on the left
commentstyle = \color{green}, %% set comment color
keywordstyle = \color{blue}, %% set keyword color
stringstyle = \color{red}, %% set string color
rulecolor = \color{black}, %% set frame color to avoid being affected by text color
basicstyle = \small \ttfamily , %% set listing font and size
breaklines = true, %% enable line breaking
numberstyle = \tiny,
}



\usetikzlibrary{calc,patterns,angles,quotes}
\usepackage{pgfplots}
\pgfplotsset{width=15cm,height=10cm,compat=1.9}

\pgfplotstableread{
x	y	delta-x	delta-y
-0.009	-6.3255	0.0284	0.0028
0.0611	-5.9201	0.0302	0.0019
0.2967	-5.6324	0.0375	0.0014
0.4689	-5.4092	0.0445	0.0011
0.4903	-5.2269	0.0455	0.0009
0.4822	-5.0728	0.0451	0.0008
0.5518	-4.9392	0.0485	0.0007
}{\mytable}

\begin{filecontents}{graphdata.dat}
x	y	delta-x	delta-y
-0.009	-6.3255	0.0284	0.0028
0.0611	-5.9201	0.0302	0.0019
0.2967	-5.6324	0.0375	0.0014
0.4689	-5.4092	0.0445	0.0011
0.4903	-5.2269	0.0455	0.0009
0.4822	-5.0728	0.0451	0.0008
0.5518	-4.9392	0.0485	0.0007
\end{filecontents}

\usepackage[english]{babel}
\usepackage{fancyhdr}
\usepackage{lastpage}
\pagestyle{fancy}
\fancyhf{}


\pagestyle{empty}% for cropping
\usepackage{tikz}
\usetikzlibrary{calc}
\usepackage{zref-savepos}
\usepackage{tabu}

\newcounter{NoTableEntry}
\renewcommand*{\theNoTableEntry}{NTE-\the\value{NoTableEntry}}

\newcommand*{\strike}[2]{%
  \multicolumn{1}{#1}{%
    \stepcounter{NoTableEntry}%
    \vadjust pre{\zsavepos{\theNoTableEntry t}}% top
    \vadjust{\zsavepos{\theNoTableEntry b}}% bottom
    \zsavepos{\theNoTableEntry l}% left
    \hspace{0pt plus 1filll}%
    #2% content
    \hspace{0pt plus 1filll}%
    \zsavepos{\theNoTableEntry r}% right
    \tikz[overlay]{%
      \draw
        let
          \n{llx}={\zposx{\theNoTableEntry l}sp-\zposx{\theNoTableEntry r}sp-\tabcolsep},
          \n{urx}={\tabcolsep},
          \n{lly}={\zposy{\theNoTableEntry b}sp-\zposy{\theNoTableEntry r}sp},
          \n{ury}={\zposy{\theNoTableEntry t}sp-\zposy{\theNoTableEntry r}sp}
        in
        (\n{llx}, \n{lly}) -- (\n{urx}, \n{ury})
      ;
    }% 
  }%
}

\usepackage{tikz}
\usetikzlibrary{calc}

\tikzset{
    right angle quadrant/.code={
        \pgfmathsetmacro\quadranta{{1,1,-1,-1}[#1-1]}     % Arrays for selecting quadrant
        \pgfmathsetmacro\quadrantb{{1,-1,-1,1}[#1-1]}},
    right angle quadrant=1, % Make sure it is set, even if not called explicitly
    right angle length/.code={\def\rightanglelength{#1}},   % Length of symbol
    right angle length=2ex, % Make sure it is set...
    right angle symbol/.style n args={3}{
        insert path={
            let \p0 = ($(#1)!(#3)!(#2)$) in     % Intersection
                let \p1 = ($(\p0)!\quadranta*\rightanglelength!(#3)$), % Point on base line
                \p2 = ($(\p0)!\quadrantb*\rightanglelength!(#2)$) in % Point on perpendicular line
                let \p3 = ($(\p1)+(\p2)-(\p0)$) in  % Corner point of symbol
            (\p1) -- (\p3) -- (\p2)
        }
    }
}
\pagestyle{fancy}
\fancyhf{}
\title{SPH4U1 Grade 12 Physics Laboratory 3:\\Investigation of Magnitude of Effect Radius Has on Gravity}
\author{Andi Tse}
\date{June 6, 2024}
\rhead{Instructor: Dr. Dykshoorn}
\lhead{Andi Tse}
\chead{Laboratory 3}
\rfoot{SPH4U1}
\lfoot{Grade 12 Physics}
\cfoot{Page \thepage \hspace{1pt} of \pageref{LastPage}}
\pagestyle{fancy}

 


\begin{document}

\maketitle

The following diagram illustrate the deviation of the path of an object under the gravitational attraction

\begin{center}
    \includegraphics[width=7cm]{High School Courses/SPH4U1/Gravitational Deviation Diagram.png}
\end{center}

To analyze the effect of radial distance have on the magnitude of gravitational force, we construct the following diagram:

\begin{center}
    \includegraphics[width=7cm]{High School Courses/SPH4U1/Gravitational Deviation Diagram2.png}
\end{center}

Newton proved in his book of \textit{Mathematical Principles of Natural Philosophy} that \[|\Vec{F}_\text{gravitational}|\propto \frac{QR}{(SP\times QT)^2}\tag{1}\]
We would assume that 
\[|\Vec{F}_\text{gravitational}|\propto SP^n\tag{2}\]
and we will attempt to find $n$. Combining proportionality 1 and 2 gives us:
\[\frac{QR}{(SP\times QT)^2}\propto SP^n \tag{*}\]


Based on the data we obtained, 
% Please add the following required packages to your document preamble:
% \usepackage{multirow}
\begin{table}[h]
\centering
\begin{tabular}{|c|ccc|l}
\cline{1-4}
\multirow{2}{*}{Set \#} & \multicolumn{3}{c|}{Length (cm)}                            &  \\ \cline{2-4}
                        & \multicolumn{1}{c|}{SP $(\pm 0.1)$}   & \multicolumn{1}{c|}{QR $(\pm 0.1)$}  & QT  $(\pm 0.1)$  &  \\ \cline{1-4}
1                       & \multicolumn{1}{c|}{20.0} & \multicolumn{1}{c|}{2.1} & 10.0 &  \\ \cline{1-4}
2                       & \multicolumn{1}{c|}{23.7} & \multicolumn{1}{c|}{2.5} & 10.2 &  \\ \cline{1-4}
3                       & \multicolumn{1}{c|}{32.7} & \multicolumn{1}{c|}{2.0} & 9.8  &  \\ \cline{1-4}
4                       & \multicolumn{1}{c|}{39.0} & \multicolumn{1}{c|}{1.4} & 9.0  &  \\ \cline{1-4}
5                       & \multicolumn{1}{c|}{39.3} & \multicolumn{1}{c|}{1.2} & 7.0  &  \\ \cline{1-4}
6                       & \multicolumn{1}{c|}{35.4} & \multicolumn{1}{c|}{1.1} & 7.0  &  \\ \cline{1-4}
7                       & \multicolumn{1}{c|}{28.6} & \multicolumn{1}{c|}{2.1} & 8.0  &  \\ \cline{1-4}
\end{tabular}
\end{table}
\\
Uncertainties would follow the formula:

\begin{align*}
    \delta\left(\ln\left(\frac{QR}{(SP\times QT)^2}\right)\right)&=\frac{\delta QR}{QR}+\frac{2\delta SP}{SP}+\frac{2\delta QT}{QT}\\
    \delta \left(\ln\left(SP\right)\right)&=\frac{\delta SP}{SP}
\end{align*}

Taking the first set of data as example, 
\begin{align*}
    \delta\left(\ln\left(\frac{QR}{(SP\times QT)^2}\right)\right)&=\frac{\delta QR}{QR}+\frac{2\delta SP}{SP}+\frac{2\delta QT}{QT}\\
    &=\frac{0.1}{2.1}+\frac{2\times 0.1}{20.0}+\frac{2\times 0.1}{10.0}\\
    &=0.01\\
    \delta \left(\ln\left(SP\right)\right)&=\frac{\delta SP}{SP}\\
    &=\frac{0.1}{20.0}\\
    &=0.08
\end{align*}
But we will promote both uncertainty to 0.1 to meet the significant figures.

Processed best data would be
\begin{table}[h]
\centering
\begin{tabular}{|c|c|c|c|}
\hline
$\ln\left(\frac{QR}{(SP\times QT)^2}\right)$ (1) & $\delta\left(\ln\left(\frac{QR}{(SP\times QT)^2}\right)\right)$(1)& $\ln(SP)$(1) & $\delta \ln(SP)$(1) \\ \hline
-9.9                              & 0.1                                     & 3.0    & 0.1          \\ \hline
-10.1                             & 0.1                                     & 3.2    & 0.1          \\ \hline
-10.8                             & 0.1                                     & 3.5    & 0.1          \\ \hline
-11.4                             & 0.1                                     & 3.7    & 0.1          \\ \hline
-11.1                             & 0.1                                     & 3.7    & 0.1          \\ \hline
-10.9                             & 0.1                                     & 3.6    & 0.1          \\ \hline
-10.1                             & 0.1                                     & 3.4    & 0.1          \\ \hline
\end{tabular}
\end{table}


The final graph is as follows:
\begin{center}
    \includegraphics[width=15cm]{High School Courses/SPH4U1/ln(SP) and ln(QR_(SP_QT)^2).png}
\end{center}
The $y$-axis is $\ln\left(\frac{QR}{(SP\times QT)^2}\right)$ and the $x$-axis is $\ln(SP)$

Reading off the graph using the points $(3.0,-9.7),(3.6,-11)$, the slope has no unit and the value would be 
\begin{align*}
    n&=\frac{-11+9.7}{3.6-3.0}=-2.2
\end{align*}

Which is off by
\[\%difference=\left|\frac{theretical-experimental}{theoretical}\times 100\%\right|=10\%\]
From the theoretical result. 

\begin{itemize}
    \item [Q1.] Why should Q → P?\\
    It is so that the radial distance $SP$ and $SQ$ are as close to each other as possible, delivering more accurate results
\item [Q2.] Why is the line segment QR parallel to SP?
\\If the point $Q$ is approaching point $P$, the path deviation would be parallel to $SP$
\item [Q3.] Is the method general enough to work for other radial force laws?\\
Newton used Kepler's second law to convert time into measurable area, the law is not necessarily true for other radial force laws. However, our part of the work, the method finding the exponent of radius works for all radial forces as long as the radial force is measurable
\end{itemize}


\end{document}
